\section{Related Work}
\label{sec:related_work}

\paragraph{Tool-augmented LLMs.}
Tool-augmented language models extend LLMs from text generation to action generation by allowing them to call external APIs, search engines, calculators, software interfaces, and other executable resources. Early work showed that models can learn when and how to invoke tools from demonstrations or self-supervised traces \citep{schick2023toolformer, yao2022react}, and subsequent agent benchmarks have emphasized planning, execution, and interaction in realistic environments \citep{yao2022webshop, huang2024planning}. A large body of work improves tool use by instruction tuning or supervised fine-tuning on tool-call trajectories. Gorilla \citep{patil2024gorilla}, ToolLLM \citep{qin2023toolllm}, API-Bank \citep{li-etal-2023-api}, ToolACE \citep{liu2024toolace}, and APIGen \citep{liu2024apigen} demonstrate that tool-specific training can improve schema following, argument generation, and function-call formatting. These approaches are complementary to Toollery: they improve the base model's ability to execute a chosen tool call, while Toollery focuses on reducing the candidate set that the model must inspect before making that call. This distinction is important in large tool libraries, where even a strong function-calling model may fail because the relevant schema is buried among hundreds or thousands of distractors.

\paragraph{Tool and skill retrieval and routing.}
A second line of work treats tool selection as retrieval or routing. Rather than presenting every available tool to the LLM, the system first retrieves a smaller set of candidate tools and then asks the LLM to choose among them. ToolRM \citep{agarwal2025toolrm}, MassTool \citep{lin2025masstool}, RAG-MCP \citep{gan2025rag}, and context-pruning approaches such as Less is More \citep{paramanayakam2025less} share this broad motivation: tool use improves when the model reasons over a compact and relevant action space instead of a long, noisy tool prompt. SkillRouter is especially related because it also aims to route user requests to a subset of skills before final model inference. However, SkillRouter performs compact full-text retrieval and reranking over the skill body or skill documentation. Toollery makes a different design choice. It does not retrieve directly over the full tool description as the primary semantic object; instead, it constructs a self-verified generated-query index offline and performs query-level retrieval at inference time. In this sense, Toollery is not a universal skill router and is not intended to replace SkillRouter. It can be viewed as a pre-routing candidate compression layer that supplies a smaller, semantically grounded candidate set to an existing router or to a standard function-calling LLM.

\paragraph{Query augmentation and verification.}
Toollery is also connected to query augmentation methods that generate alternative user intents, pseudo-queries, or synthetic supervision in order to improve retrieval recall. In document retrieval, generated queries can make sparse or dense indexes more robust to vocabulary mismatch; in tool selection, the analogous mismatch is between user intent and tool metadata. A tool description usually states what an API does, while a user request states the goal they want to accomplish. Directly embedding the tool description and the user request therefore asks the retriever to bridge a substantial semantic gap. Toollery addresses this gap by generating diverse natural-language queries for each tool and then filtering them through round-trip self-verification. The verification step is central: a generated query is retained only when an independent verifier can recover the originating tool from a candidate set containing distractors. The resulting verified queries are not merely paraphrases of the tool description; they are high-precision semantic anchors that have survived an explicit ambiguity check. This makes the index better suited for candidate compression, where false positives can crowd out the correct tool under a fixed top-$k$ budget.

\paragraph{RAG baselines and generic document retrieval.}
Retrieval-augmented generation systems retrieve external textual evidence and condition the LLM on the retrieved documents. Generic RAG frameworks, including systems such as RAGAnything, are designed for broad document grounding across heterogeneous corpora and modalities. Toollery differs in both the indexed unit and the downstream contract. A generic RAG baseline retrieves chunks of documentation and sends them to the LLM as evidence. Toollery indexes verified queries, uses them only to identify candidate tools, and sends the LLM compact tool specifications rather than arbitrary retrieved text. This distinction matters because the final task is not open-ended answer generation but constrained tool selection and argument generation. The retrieval layer should therefore optimize for selection quality and context compression, not for document coverage. Toollery can be implemented using standard dense retrieval infrastructure, but its contribution lies in the construction and use of a verified query index specialized for tool candidate selection.

\section{The Toollery Design}
\label{sec:method_late_draft}

\subsection{Terminology Update}
\label{subsec:terminology_update}

The revised framing of Toollery replaces earlier manual-centric terminology with retrieval-centric terminology. The term \textit{proxy query} is replaced by \textit{verified query}: each generated query is retained only after passing a round-trip verification test against distractor tools. The term \textit{Tool Manual} or \textit{Tool User Manual} is replaced by \textit{verified query index}: the offline artifact is an index of verified intent--tool pairs, not a human-facing manual or a full documentation store. Likewise, \textit{proxy-query matching} is replaced by \textit{verified query retrieval}, and \textit{manual synthesis} is replaced by \textit{verified query index construction}. These changes reflect the intended role of Toollery: it is a self-verified generated-query retrieval framework and a pre-routing candidate compression layer, rather than a full skill router or a generic documentation RAG system.

\begin{table}[t]
\centering
\small
\begin{tabular}{ll}
\hline
Earlier term & Revised term \\
\hline
proxy query & verified query \\
Tool Manual / Tool User Manual & verified query index \\
proxy-query matching & verified query retrieval \\
manual synthesis & verified query index construction \\
full tool-pool prompting & uncompressed candidate prompting \\
\hline
\end{tabular}
\caption{Terminology used in the revised Toollery narrative.}
\label{tab:terminology_update}
\end{table}

\subsection{Verified Query Retrieval}
\label{subsec:verified_query_retrieval}

After the offline construction stage, Toollery has an index $\mathcal{I}$ containing verified query--tool pairs. Each entry is a tuple $(q_i, t_i, s_i)$, where $q_i$ is a generated user-like query, $t_i \in \mathcal{T}$ is the tool that the query was generated for, and $s_i$ optionally stores the verification signal, such as the verifier's confidence or the rank of the target tool during verification. The key property of $\mathcal{I}$ is that each query has been tested for functional alignment: when the verifier receives $q_i$ together with the target tool and sampled distractors, it must select $t_i$. Toollery therefore retrieves over a space of verified intents rather than over raw tool descriptions.

At inference time, given a user request $x$, Toollery embeds $x$ and the verified queries into a shared dense space. Let $e(x)$ denote the embedding of the user request and $e(q_i)$ denote the embedding of a verified query. The retrieval score for each verified query is
\begin{equation}
    r_i = \cos(e(x), e(q_i)).
\end{equation}
Toollery then retrieves the top-$m$ verified queries by $r_i$. This stage is intentionally lightweight: it can be implemented with any standard vector index and does not require updating the backbone LLM. The important difference from direct tool retrieval is the semantic target. Direct retrieval asks whether a user request resembles a tool description, which is often written in API-centric language. Verified query retrieval asks whether a user request resembles another user-like intent that has already been validated as mapping to a tool.

Because multiple verified queries may point to the same tool, Toollery aggregates query-level evidence into tool-level scores. For a tool $t$, let $\mathcal{N}_m(t)$ be the set of retrieved verified queries among the top-$m$ whose target tool is $t$. A simple aggregation function is
\begin{equation}
    R(t \mid x) = \max_{q_i \in \mathcal{N}_m(t)} r_i,
\end{equation}
which rewards the best-matching verified intent for each tool. In practice, other aggregators can be used, such as the mean of the top few scores, a log-sum-exp score, or a weighted score that incorporates verification confidence $s_i$. We use aggregation to convert many query-level matches into a stable ranked list of tools while avoiding repeated copies of the same tool in the final context.

\subsection{Candidate Subset Construction and LLM Selection}
\label{subsec:candidate_subset}

The output of verified query retrieval is not the final tool call. Instead, it is a compressed candidate subset passed to the LLM for final selection and argument generation. Toollery ranks tools by $R(t \mid x)$ and selects the top-$k$ unique tools:
\begin{equation}
    \mathcal{C}_k(x) = \operatorname{TopK}_{t \in \mathcal{T}} R(t \mid x).
\end{equation}
For each selected tool, Toollery instantiates the compact tool specification required by the target function-calling interface, including the tool name, natural-language description, and parameter schema. The LLM then receives the original user request $x$ together with $\mathcal{C}_k(x)$ and produces the final tool call.

This two-step design preserves the role of the LLM as the final decision-maker. Toollery does not assume that nearest-neighbor retrieval alone is sufficient for tool use, especially when tools have overlapping descriptions or when the user request requires careful argument extraction. Instead, retrieval narrows the action space and the LLM resolves the remaining ambiguity. This makes Toollery compatible with native function-calling APIs, prompt-based tool calling, and external routers. When used with SkillRouter-like systems, Toollery can operate before the router as a candidate compression layer; when used without an additional router, the compressed subset can be sent directly to the backbone LLM.

The top-$k$ value controls the recall--compression trade-off. A larger $k$ increases the probability that the ground-truth tool is included, but also increases prompt length and distractor pressure. A smaller $k$ improves efficiency and reduces irrelevant schemas, but can hurt accuracy if retrieval misses the correct tool. The central empirical question is therefore whether verified query retrieval can maintain high candidate recall with small $k$ as the global tool pool grows. Our experiments evaluate this question under controlled candidate-set expansion.

\subsection{Token and Latency Analysis}
\label{subsec:token_latency}

Toollery is motivated not only by selection quality but also by the token and latency costs of uncompressed tool prompting. Let $N$ be the number of tools in the full library, and let $L_t$ be the average token length of a tool specification. A standard full-prompting approach sends approximately
\begin{equation}
    O(N \cdot L_t)
\end{equation}
tool-definition tokens to the LLM for each request, in addition to the user prompt and system instructions. As $N$ grows from tens to hundreds or thousands, the prompt can exceed the model's context window or consume most of the available budget before the model has reasoned about the user request.

Toollery changes the online LLM context from full-library prompting to top-$k$ candidate prompting. The LLM receives approximately
\begin{equation}
    O(k \cdot L_t)
\end{equation}
tool-definition tokens, where $k \ll N$ and is fixed independently of the full library size. The retrieval index still grows with the tool library, but dense vector search is substantially cheaper than long-context LLM inference and can be served using standard approximate nearest-neighbor infrastructure. The offline cost of generating and verifying queries is amortized over many future requests and only needs to be rerun for new or modified tools.

This distinction is important for latency. In full-prompting systems, both prefill cost and attention cost increase with the number of candidate tools sent to the model. In Toollery, the LLM-facing prompt length remains nearly constant once $k$ is fixed. The online pipeline adds an embedding and retrieval step, but this step replaces a much larger LLM context expansion. Toollery should therefore be evaluated as a selection-quality and compression mechanism: it aims to preserve or improve final tool-call accuracy while reducing the number of tool tokens and the amount of irrelevant schema information presented to the LLM.

\section{Discussion}
\label{sec:discussion_draft}

The central claim of Toollery is deliberately narrow. Toollery is not a universal skill router, a replacement for trained function-calling models, or a general-purpose RAG system over tool documentation. It is a self-verified generated-query retrieval framework for compressing the candidate tool set before LLM selection. This framing clarifies both where Toollery helps and where it should be combined with other components.

The experiments support two main observations. First, tool selection degrades as the candidate set grows, even when the ground-truth tool and user request remain unchanged. This suggests that large tool libraries introduce a search-space problem in addition to the familiar challenges of reasoning and argument generation. Second, verified query retrieval can reduce this search-space pressure by presenting the LLM with a compact set of semantically plausible tools. The improvement comes from changing what the model is asked to do: instead of scanning a large library of heterogeneous schemas, the model performs final selection within a small candidate subset.

This design also explains why Toollery can help both small and large models. Smaller models often have limited context processing ability and are more vulnerable to long prompts containing many irrelevant schemas. Larger models may tolerate longer contexts, but they still pay increased token and latency costs and can be distracted by semantically similar tools. Toollery reduces both types of burden. It compresses the prompt while preserving the final LLM step needed for schema-aware selection and argument filling.

The comparison with SkillRouter and RAGAnything highlights the intended scope. SkillRouter retrieves and reranks over skill bodies; Toollery retrieves over verified generated queries that are used only to produce a candidate set. RAGAnything and related generic RAG systems retrieve document content for grounding; Toollery retrieves intent anchors and ultimately sends compact tool specifications to the model. These systems can be complementary. A deployment could first use Toollery to reduce a large tool library to a manageable subset, then apply a router, reranker, or function-calling model within that subset.

\section{Limitations}
\label{sec:limitations_draft}

Toollery depends on the quality of the verified query index. If the teacher model fails to generate queries that cover realistic user intents, retrieval recall will be limited no matter how strong the final LLM is. Conversely, if verification is too weak or the distractor set is not challenging enough, ambiguous verified queries may enter the index and produce false-positive candidates. The verification procedure therefore improves precision but does not provide a formal guarantee of correctness.

The current framework also introduces offline maintenance cost. New tools require query generation, verification, embedding, and index updates. This cost is amortized across future calls, but it matters for rapidly changing tool ecosystems. Incremental index construction and targeted re-verification are natural extensions for production settings where tools are frequently added, removed, or modified.

Another limitation is that Toollery primarily addresses single-step candidate selection. It does not by itself solve multi-tool planning, tool sequencing, execution monitoring, or recovery from failed calls. In multi-step agents, Toollery can be applied at each decision point, but the broader planning policy remains outside the scope of the method. Similarly, Toollery compresses the candidate set but does not guarantee correct argument values; the final LLM must still understand the selected schema and fill parameters accurately.

Finally, Toollery inherits the limitations of embedding-based retrieval. Dense retrieval can miss rare intents, domain-specific terminology, or requests that require compositional reasoning across multiple tools. Hybrid retrieval, stronger query aggregation, learned reranking, or domain-adaptive encoders may further improve candidate recall, but these additions should be evaluated against their added complexity and latency.

\section{Conclusion}
\label{sec:conclusion_draft}

We presented Toollery, a self-verified generated-query retrieval framework for scalable tool selection. The key idea is to move part of the semantic grounding work offline: for each tool, Toollery generates diverse user-like queries and retains only those that pass round-trip verification against distractor tools. At inference time, the user request is matched against this verified query index, producing a compact candidate subset that is then passed to the LLM for final tool selection and argument generation.

This framing positions Toollery as a pre-routing candidate compression layer rather than a universal router or a generic RAG system. Its goal is to improve selection quality while reducing token and latency costs under large tool pools. By retrieving over verified intents and sending only top-$k$ compact tool specifications to the LLM, Toollery reduces the amount of irrelevant schema information in the prompt and helps existing function-calling pipelines operate more reliably as tool libraries grow.
